\begin{proof}[Proof of \cref{obj:thm:tv-envelope-design}]
Let
\[
c:=\frac{\sigma_0^2}{n}.
\]
The envelope-scale assumption gives \(c\ge 0\). We use this nonnegativity twice below, when pulling the scalar \(c\) through an infimum. % lean: div_nonneg

\begin{enumerate}
\item
First, \(W_\beta(p)\) is nonempty. Indeed, the order hypothesis gives \(\beta\ge 1\), the theorem assumes \(k\ge \beta\), and \(p\in S_{k,q}\). Hence \cref{obj:lem:unbiased-weight-set-nonempty} applies and gives a weight vector \(w\in W_\beta(p)\). % lean: unbiased_weight_set_nonempty

\item
Fix now any \(w\in W_\beta(p)\). By linearity of design expectation and by the defining mean-curve restriction in \(\mathcal P_\beta\),
\[
\mathbb E_\pi\!\left[\sum_{j=0}^k w_j\bar Y_j\right]
=
\sum_{j=0}^k w_j\,m_P(p_j).
\]
Since \(p\in S_{k,q}\), every node \(p_j\) lies in \([0,1]\), so the degree-\(\beta\) polynomial condition in \(\mathcal P_\beta\) gives
\[
m_P(p_j)=\sum_{\ell=0}^{\beta} a_{P,\ell}p_j^\ell .
\]
Substituting and interchanging the two finite sums,
\[
\sum_{j=0}^k w_jm_P(p_j)
=
\sum_{\ell=0}^{\beta}
a_{P,\ell}\sum_{j=0}^k w_jp_j^\ell .
\]
The moment equations defining \(W_\beta(p)\) give
\[
\sum_{j=0}^k w_jp_j^0=0,
\qquad
\sum_{j=0}^k w_jp_j^\ell=1
\quad(1\le \ell\le \beta).
\]
Thus the last display equals \(\sum_{\ell=1}^{\beta}a_{P,\ell}\). By \cref{obj:prop:rollout-polynomial-identity}, this is \(m_P(1)-m_P(0)\), proving the unbiasedness identity. % lean: rollout_polynomial_identity

\item
For the variance statement, apply \cref{obj:lem:variance-envelope-sharpness} with \(X_j=\bar Y_j\). Its hypotheses are exactly the nonnegative-scale assumption and the round-mean variance-envelope component of \(\mathcal P_\beta\). Therefore
\[
\operatorname{Var}_\pi\!\left(\sum_{j=0}^k w_j\bar Y_j\right)
\le
\frac{\sigma_0^2}{n}\left(\sum_{j=0}^k |w_j|\right)^2,
\]
and the same lemma supplies a positive semidefinite matrix \(\Gamma\) with
\[
\Gamma_{jj}\le \frac{\sigma_0^2}{n}
\quad\text{for all }j,
\qquad
\sum_{i=0}^k\sum_{j=0}^k w_i\Gamma_{ij}w_j
=
\frac{\sigma_0^2}{n}\left(\sum_{j=0}^k |w_j|\right)^2 .
\] % lean: variance_envelope_sharpness

\item
For the fixed schedule \(p\), the defining set on the left-hand side of the first envelope identity is
\[
\left\{c\,u:\ 
u=\left(\sum_{j=0}^k |w_j|\right)^2
\text{ for some }w\in W_\beta(p)
\right\}.
\]
Since \(c\ge0\), the infimum of this scalar multiple is \(c\) times the infimum of the unscaled set. By the definition of \(A_\beta(p)\) in \cref{obj:def:amplification-criterion}, this gives
\[
\inf\left\{
v:\ \exists w\in W_\beta(p),\
v=\frac{\sigma_0^2}{n}\left(\sum_{j=0}^k |w_j|\right)^2
\right\}
=
\frac{\sigma_0^2}{n}A_\beta(p).
\] % lean: Real.sInf_smul_of_nonneg

\item
The same scalar-infimum argument applies after also varying the schedule. Namely,
\[
\left\{
v:\ \exists p'\in S_{k,q},\
v=c\,A_\beta(p')
\right\}
=
c\left\{
u:\ \exists p'\in S_{k,q},\
u=A_\beta(p')
\right\}.
\]
Using \(c\ge0\) and the definition of \(M_{\beta,k,q}\) in \cref{obj:def:amplification-criterion}, we obtain
\[
\inf\left\{
v:\ \exists p'\in S_{k,q},\
v=\frac{\sigma_0^2}{n}A_\beta(p')
\right\}
=
\frac{\sigma_0^2}{n}M_{\beta,k,q}.
\] % lean: Real.sInf_smul_of_nonneg

\item
Finally, suppose \(p^\star\in S_{k,q}\) and
\[
A_\beta(p^\star)=M_{\beta,k,q}.
\]
For any \(p'\in S_{k,q}\), the definition of \(M_{\beta,k,q}\) as the infimum over budgeted schedules gives
\[
M_{\beta,k,q}\le A_\beta(p').
\]
Hence \(A_\beta(p^\star)\le A_\beta(p')\). Multiplying by the nonnegative scalar \(c=\sigma_0^2/n\) yields
\[
\frac{\sigma_0^2}{n}A_\beta(p^\star)
\le
\frac{\sigma_0^2}{n}A_\beta(p'),
\]
as required. % lean: minimaxAmplification_le_of_budgeted
\end{enumerate}
\end{proof}
